\documentclass[12pt]{article}

\usepackage{sbc-template}
\usepackage{graphicx,url}
\usepackage[utf8]{inputenc}
\usepackage{float}

\sloppy

\title{Analise de Colaboracao no GitHub usando Grafos}

\author{Arthur\inst{1}, Amanda\inst{1}, Pedro\inst{1}}

\address{Instituto de Ciencias Exatas e Informatica\\
Pontificia Universidade Catolica de Minas Gerais (PUC Minas)
}

\begin{document}

\maketitle

\begin{abstract}
This report describes a Java tool for modeling collaboration in an open-source
GitHub repository as directed weighted graphs. Users are represented as
vertices, and technical interactions such as comments, issue closures, reviews
and merges are represented as edges. The implementation provides a custom graph
API with adjacency matrix and adjacency list representations, GitHub data
extraction, Gephi export and manual network metrics.
\end{abstract}

\begin{resumo}
Este relatorio descreve uma ferramenta em Java para modelar a colaboracao em
um repositorio publico do GitHub usando grafos direcionados e ponderados. Os
usuarios sao representados como vertices, e interacoes tecnicas como
comentarios, fechamento de issues, revisoes e merges sao representadas como
arestas. A implementacao fornece uma API propria de grafos com matriz e lista
de adjacencia, extracao de dados do GitHub, exportacao para Gephi e calculo
manual de metricas de redes.
\end{resumo}

\section{Introducao}

O objetivo do trabalho e desenvolver uma ferramenta computacional baseada em
Teoria dos Grafos para estudar interacoes de colaboradores em um projeto real
de codigo aberto. A ferramenta foi implementada em Java 21, com Maven, e evita
bibliotecas prontas de grafos. Essa decisao foi tomada para atender diretamente
ao requisito de implementar a estrutura de grafos e seus algoritmos de forma
propria.

O repositorio escolhido foi \texttt{encode/httpx}, um cliente HTTP moderno para
Python. A escolha se justifica por ser um projeto publico, conhecido, com mais
de 5.000 estrelas e com historico de issues, pull requests, comentarios e
revisoes. Esse conjunto de eventos permite construir uma rede de colaboracao
com diferentes tipos de relacionamento entre usuarios.

\section{Responsabilidades}

O trabalho foi organizado entre Arthur, Amanda e Pedro, com divisao de tarefas
para cobrir modelagem, implementacao, validacao e escrita do relatorio:

\begin{itemize}
  \item Arthur ficou responsavel pela organizacao da solucao, implementacao da
  API de grafos em Java, adaptacao da arquitetura, testes automatizados e
  documentacao tecnica.
  \item Amanda ficou responsavel pela revisao da modelagem do problema,
  definicao das interacoes consideradas, verificacao dos pesos do grafo
  integrado e apoio na escrita do relatorio.
  \item Pedro ficou responsavel pelo apoio na mineracao dos dados do GitHub,
  validacao dos arquivos exportados para o Gephi, revisao das metricas de rede
  e organizacao dos resultados para a entrega final.
\end{itemize}

\section{Modelagem do Grafo}

Cada usuario do GitHub e modelado como um vertice. Cada interacao entre dois
usuarios e modelada como uma aresta direcionada. O grafo e simples: nao permite
lacos nem multiplas arestas. Quando ocorre uma nova interacao entre o mesmo par
ordenado de usuarios, a aresta existente tem seu peso acumulado.

Foram definidos tres grafos separados:

\begin{itemize}
  \item grafo de comentarios em issues e pull requests;
  \item grafo de fechamento de issue por outro usuario;
  \item grafo de revisoes, aprovacoes e merges de pull requests.
\end{itemize}

Alem desses grafos separados, tambem foi definido um grafo integrado
ponderado. Nele, os pesos representam a intensidade da relacao tecnica.

\begin{table}[H]
\centering
\caption{Pesos usados no grafo integrado}
\begin{tabular}{lr}
\hline
Interacao & Peso \\
\hline
Comentario em issue ou pull request & 2 \\
Abertura de issue comentada por outro usuario & 3 \\
Fechamento de issue por outro usuario & 3 \\
Revisao ou aprovacao de pull request & 4 \\
Merge de pull request & 5 \\
\hline
\end{tabular}
\end{table}

O fechamento de issue foi mantido no grafo integrado com peso 3 por representar
uma acao de triagem relevante, de intensidade intermediaria. Essa adaptacao e
coerente com a permissao do enunciado para ajustar os pesos desde que as
interacoes tecnicas mais fortes tenham maior peso que interacoes leves.

\section{Arquitetura da Ferramenta}

A arquitetura foi organizada em camadas. O pacote \texttt{core.graph} contem a
API obrigatoria do enunciado, incluindo \texttt{AbstractGraph},
\texttt{AdjacencyMatrixGraph} e \texttt{AdjacencyListGraph}. Essas classes
mantem os nomes em ingles porque o enunciado os solicita literalmente. Para
facilitar o estudo, as implementacoes principais tambem possuem nomes em
portugues, como \texttt{GrafoMatrizAdjacencia} e
\texttt{GrafoListaAdjacencia}.

Os demais pacotes foram padronizados em portugues:

\begin{itemize}
  \item \texttt{aplicacao}: comandos de terminal e demonstracao da API;
  \item \texttt{extrator}: configuracao e cliente GraphQL do GitHub;
  \item \texttt{leitura}: conversao do JSON do GitHub em interacoes e grafos;
  \item \texttt{arquivos}: exportacao e leitura de arquivos CSV do Gephi;
  \item \texttt{estatisticas}: calculo manual das metricas de rede;
  \item \texttt{servicos}: construcao de grafos a partir de interacoes.
\end{itemize}

A aplicacao principal aceita os comandos \texttt{fetch}, \texttt{build} e
\texttt{analyze}. Esses nomes foram mantidos como comandos de terminal por
serem curtos e comuns em ferramentas de linha de comando, mas a implementacao
interna foi nomeada em portugues.

\section{API de Grafos}

A classe abstrata \texttt{AbstractGraph} define os atributos comuns e os
metodos obrigatorios, como quantidade de vertices, quantidade de arestas,
verificacao de aresta, graus de entrada e saida, pesos, conectividade, grafo
vazio, grafo completo e exportacao para o Gephi.

Foram implementadas duas representacoes:

\begin{itemize}
  \item matriz de adjacencia, adequada para grafos densos;
  \item lista de adjacencia, adequada para grafos esparsos.
\end{itemize}

As operacoes validam indices invalidos e operacoes inconsistentes. Lacos sao
rejeitados, e arestas repetidas nao sao duplicadas. Quando uma interacao ja
existe, o peso e somado na mesma aresta.

\section{Coleta de Dados}

A coleta de dados foi implementada no pacote \texttt{extrator}, usando
\texttt{HttpClient} do Java e a API GraphQL do GitHub. A consulta busca issues
fechadas, pull requests recentes, comentarios, revisoes, usuario responsavel
pelo merge e eventos de fechamento de issue.

Para executar a coleta real, e necessario configurar um arquivo \texttt{.env}
com \texttt{GITHUB\_TOKEN}, \texttt{GITHUB\_REPO\_OWNER},
\texttt{GITHUB\_REPO\_NAME}, \texttt{MAX\_ISSUES} e \texttt{MAX\_PULLS}. No
ambiente local usado nesta revisao ainda nao havia token nem dados reais
gerados. Portanto, a ferramenta esta pronta para a coleta, mas a execucao real
deve ser feita antes da entrega final.

\section{Metricas de Analise}

As metricas foram implementadas manualmente, sem NetworkX ou bibliotecas
prontas de grafos. A estrategia \texttt{EstatisticasGrafoManual} calcula:

\begin{itemize}
  \item centralidade de grau, entrada e saida;
  \item centralidade de intermediacao;
  \item centralidade de proximidade;
  \item PageRank;
  \item centralidade de autovetor;
  \item densidade;
  \item coeficiente de aglomeracao;
  \item assortatividade;
  \item comunidades;
  \item modularidade;
  \item vertices de ponte entre comunidades.
\end{itemize}

As metricas dos vertices sao exportadas para \texttt{nodes.csv}, e as metricas
gerais do grafo sao exportadas para \texttt{graph.json}. Esses arquivos podem
ser usados no Gephi e no texto final do relatorio.

\section{Testes}

O projeto possui testes automatizados em JUnit. Os testes cobrem a API
obrigatoria de grafos, as duas representacoes, construcao de grafos de
colaboracao, leitura de JSON de exemplo, exportacao e recarga de CSVs do Gephi
e metricas principais.

Na ultima execucao local, a suite apresentou:

\begin{verbatim}
Tests run: 41, Failures: 0, Errors: 0, Skipped: 0
BUILD SUCCESS
\end{verbatim}

\section{Resultados Esperados}

Apos configurar o token do GitHub e executar a mineracao real, espera-se gerar
os grafos separados e o grafo integrado para o repositorio \texttt{encode/httpx}.
Com os arquivos CSV e JSON resultantes, o relatorio final deve incluir tabelas
com os colaboradores mais centrais, densidade da rede, modularidade,
comunidades encontradas e interpretacao dos vertices que atuam como pontes.

Esta secao deve ser substituida pelos resultados reais antes da entrega final.
Ela foi mantida neste documento para deixar explicito o que ainda precisa ser
executado para cumprir integralmente a Etapa 3 do enunciado.

\section{Conclusao}

A ferramenta implementada atende a base estrutural solicitada: possui API
propria de grafos, duas representacoes, validacoes, exportacao para Gephi,
mineracao GitHub preparada, construcao de grafos por tipo de interacao e
metricas manuais. A padronizacao em portugues melhora a legibilidade para a
apresentacao, enquanto os nomes em ingles exigidos pelo enunciado foram
preservados por compatibilidade.

Para a entrega final, ainda e necessario executar a mineracao real, preencher
os resultados obtidos, gerar o PDF do relatorio pelo LaTeX, preparar a
apresentacao oral e, se exigido pela disciplina, gravar o video demonstrativo.

\bibliographystyle{sbc}
\bibliography{referencias}

\end{document}
